% =====================================================================
%  統計基礎 0. 準備：確率と統計の基本（範囲表外・2級相当）
%  構成は Docs/SectionPlan.md にしたがう。
% =====================================================================
\chapter{確率と統計の基本}
\label{chap:stat-preparation}

\begin{chaptergoal}
  % TODO: この章の目標
\end{chaptergoal}

% 範囲表には載っていない（2級相当）。範囲表の注記「2級・3級などの内容も出題される」に対応する。

% ---------------------------------------------------------------------
\section{データの要約}
\label{sec:stat-summary}

% TODO: Σ（シグマ）記号の使い方
% TODO: 平均・分散・標準偏差
% TODO: 共分散と相関係数
% TODO: 度数分布表、ヒストグラム、箱ひげ図

% ---------------------------------------------------------------------
\section{確率の基本}
\label{sec:stat-probability-basics}

% TODO: 試行と事象、和事象・積事象・余事象
% TODO: 確率の加法定理
% TODO: 条件付き確率と乗法定理
% TODO: 事象の独立
% TODO: ベイズの定理（簡単な例：検査の陽性と病気の確率）

% ---------------------------------------------------------------------
\section{確率変数と期待値}
\label{sec:stat-random-variable}

% TODO: 離散型の確率変数と確率分布
% TODO: 連続型の確率変数と確率密度関数（「面積が確率」というイメージ）
% TODO: 累積分布関数
% TODO: 期待値と分散の定義と性質（$E[aX+b]$、$V[aX+b]$ など）
% TODO: 2つの確率変数：同時分布、周辺分布、共分散、独立
% TODO: 和の期待値と分散（独立なときの分散の足し算）

% ---------------------------------------------------------------------
\section{基本的な確率分布}
\label{sec:stat-basic-distributions}

% TODO: 離散型：ベルヌーイ分布、二項分布、ポアソン分布、幾何分布
% TODO: 連続型：一様分布、指数分布、正規分布
% TODO: 標準化と標準正規分布表の使い方

% ---------------------------------------------------------------------
\section{推定と検定の基本}
\label{sec:stat-estimation-testing-basics}

% TODO: 母集団と標本、統計量
% TODO: 標本平均の期待値と分散
% TODO: 推定量と不偏性（不偏分散で $n-1$ で割る理由）
% TODO: 仮説検定の流れ：帰無仮説と対立仮説、有意水準、棄却域、p 値
% TODO: 第1種の過誤・第2種の過誤、検出力
% TODO: 片側検定と両側検定

\begin{chaptersummary}
  % TODO: この章のまとめ
\end{chaptersummary}
